\section{FACTS}
\subsection{Proofs for tail bound probabilistic analysis}
\label{app:proofs}

We now complete the proofs of lemmas and theorems in \cref{sec:tailbounds}.

We start with the following standard way to approximate numbers near
1 with exponentials.
\begin{lemma}\label{lem:expapprox}
    For any real constant $\alpha>0$, and any
    real $x$ with $0<x\le\alpha$, we have
    \[
    \exp\left(-\tfrac{1}{\alpha}\ln\tfrac{1}{1-\alpha}\cdot{}x\right)
    \le 1-x < \exp(-x).
    \]
\end{lemma}

We also re-state this straightforward consequence of the Hoeffding/Chernoff bound on the sum of random variables:
\begin{lemma}\label{lem:hoeff}
    Let $X_1,\ldots,X_n$ be independent Poisson trials, and write $Y=\sum_i X_i$ for their sum. If $\Expect[Y] = \mu$, then for any $\delta>0$, each of
    \(\Pr(Y \ge \mu + \delta)\) and  \(\Pr(Y \le \mu - \delta)\)
    are at most \(\exp(-2\delta^2/n)\).
\end{lemma}

We now recall and prove the building-block lemmas from \cref{sec:tailbounds}.

\lemmalowerp*

\begin{proof}
    From \eqref{eqn:pw}, we know this probability is exactly
    $p_w = 1 - \frac{\falling{(s-u)}{w}}{\falling{s}{w}},$
    where $w = v-\tau$ is the number of unfilled slots remaining.
    Using \cref{lem:expapprox} we have
    $\frac{\falling{(s-u)}{w}}{\falling{s}{w}}
    \le \left(1-\frac{u}{s}\right)^w
    \le \exp(-uw/s)$,
    which means that
    \[p_w \ge 1 - \exp\left(-\tfrac{u(v-\tau)}{s}\right)
        = 1 - \exp\left(-\tfrac{u\tau}{s}\cdot{}\left(\tfrac{v}{\tau}-1\right)\right).
    \]
    Applying the two lower bounds on $\frac{u\tau}{s}$ and $\tfrac{v}{\tau}$
    from the lemma statement yields the claimed result.
\end{proof}

\lemmaupperp*

\begin{proof}
    Using again \eqref{eqn:pw}, the probability is exactly
    $p_w = 1 - \frac{\falling{(s-u)}{w}}{\falling{s}{w}},$
    where again $w\le v$ is the number of unfilled slots for
    item $x$. Then
    \[
        \tfrac{\falling{(s-u)}{w}}{\falling{s}{w}}
        \ge \tfrac{\falling{(s-u)}{v}}{\falling{s}{v}}
        \ge \left(\tfrac{s-u-v+1}{s-v+1}\right)^v
        > \left(1 - \tfrac{u}{s-v}\right)^v.
    \]
    Using upper bounds on $\frac{v}{s}$ and $u$ from the lemma statement, we have
    \[p_w
        < 1 - \left(1 - \tfrac{u}{s-v}\right)^v
        \le 1 - \left(1 - \CzeroUpper{}\tfrac{1}{v}\right)^v.\]
    Finally, the lower bound on $v$ from the lemma statement shows $\CzeroUpper{}/v \le \CboxUpper{}$, and so we can finally use the lower exponential bound of \cref{lem:expapprox} to obtain the stated result.
\end{proof}

\lemmauppertau*

\begin{proof}
    The tipping point $\tau$ is the expected number of slots filled in the table
    if $t$ of the $m$ total calls to \Increment{} were actually called on this
    particular item.
    
    We can divide the calls to \Increment{} into two groups: the $t$ calls for
    item $x$, and the $m-t$ calls for other items. The expected number of slots
    within $x$'s item set filled by the first group is at most
    $\Cphigh t$, from \cref{lem:upperp}.
    
    For the second group, these calls to
    \Increment{} on unrelated items are distributed uniformly at random among
    all table indices, and so their expected fraction within this item set
    is the same as their overall fraction in the table. Therefore, the expected
    number of slots filled by calls to \Increment{} on other items is at most
    \[\frac{(m-t)v}{s} < \frac{mv}{s} 
        \le \frac{\CmaxvOvert}{\CminsOvern} t.\]
    
    By linearity of expectation, we can sum these two to obtain an upper bound
    on the total expected tipping point as given in the lemma statement.
\end{proof}

Now we can proceed to the proofs of the main theorems on the accuracy of the CCBF.

\theoremfp*

\begin{proof}
    Let $\tau_t$ be the tipping point for any actual number $m\le n$ of total set
    bits in the table $T$ and for the given threshold $t$.
    And consider random variables $X_1,\ldots,X_v$ for the $v$ slots assigned to
    item $x$, where each $X_1$ is 0 or 1 depending on whether the corresponding
    slot in table $T$ is 0 or 1. We want to know the probability that the sum
    of the $X_i$'s is at least $\tau_t$, which is what would cause $\TestCount(x,t)$
    to produce a false positive.
    
    Let $k = t - \Cminfp\sqrt{\lambda t}$ be the actual number of calls to
    \Increment{} on item $x$, and write $\tau_k$ for the tipping point at threshold
    $k$. By definition and the exact calculations for $\tau_k$ outlined earlier,
    we know that $\Expect[\sum X_i] = \tau_k$.
    
    The difference between these two tipping points $\tau_t - \tau_k$ is
    the expected number of extra slots filled by $t-k$ calls to \Increment{},
    which from \cref{lem:lowerp} is at least 
    \[\Cplow(t-k) = \Cplow{}\cdot{}\Cminfp{}\sqrt{\lambda t} \ge \sqrt{\tfrac{\lambda v}{2}},\]
    where in the last step we used the upper bound on $v$ from the assumptions
    of the theorem.
    
    The variables $X_i$ are not independent, but
    they are \emph{negatively correlated}, meaning that the whenever one slot is filled,
    it only decreases the likelihood that another is filled; intuitively, this is
    because there are now fewer chances to fill the other slot.
    Therefore we can apply the Hoeffding bound in this direction
    (\cref{lem:hoeff}) to say that
    \[Pr\left(\sum X_i \ge \tau_k + \sqrt{\tfrac{\lambda v}{2}}\right)
    \le \exp(-\lambda),\]
    as required.
\end{proof}

\theoremfn*

\begin{proof}
   Writing $k$ for the actual number of complaints given in
   \eqref{eqn:fn}, we need a tail bound on the probability that,
   after $k$ calls to $\Increment{}$ on the same item $x$, there are still
   fewer than $\Cfour{}t$ slots of $x$'s item set filled in, where the latter
   constant comes from applying the upper bound on the tipping point from
   \cref{lem:uppertau}.
   
   For this, we need a lower bound on the \emph{expected} number of
   bits set after $k$ calls to \Increment{} on item $x$; from \cref{lem:lowerp}
   this is at least $\Cplow{}k$.
   
   Now we can apply the Hoeffding bound (\cref{lem:hoeff}),
   with $\mu = \Cplow{}k$ and $\mu+\delta=\Cfour{}t\ge\tau$
   to see that the probability that less than $\tau$ bits of $x$'s user set
   are flipped is at most
   \begin{align*}
    &\exp\left(-2(\Cfour{}t-\Cplow{}k)^2/k\right) \\
    \le& \exp\left(-2\left(0.38\lambda + 0.66\sqrt{\lambda t}\right)^2/k\right)\\
    \le& \exp\left(-\frac{.28\lambda^2 + \lambda\sqrt{\lambda t} + .87\lambda t}%
        {.4\lambda + .7\sqrt{\lambda t} + 1.1t}\right) \\
    \le& \exp(-.7\lambda)  \le 2^{-\lambda}.
   \end{align*}
\end{proof}